% GENERATED FILE. Edit canonical lessons, then run npm run generate:books.
% canonical-source-hash: fnv1a64:a7893c5e35328712
% canonical-lessons: MR-C83-uchalne, MR-C83-odhne, MR-C83-dhakalne, MR-C83-lapavne, MR-C83-jinkne

\chapter{To lift, To pull, To push, To hide, To win}
\label{ch:mr-to-lift-to-pull-to-push-to-hide-to-win}
\hlchaptermodality{83}

\begin{chapteropening}
\textbf{By the end of this chapter:} \emph{I can say to lift, to pull, to push, to hide and to win.}

Complete the last lesson of chapter 83: I can say to lift, to pull, to push, to hide and to win.
\end{chapteropening}

\section[\texorpdfstring{\mr{उचलणे} (uchalṇe)}{उचलणे (uchalṇe)}]{\mr{उचलणे} (uchalṇe) — to lift}

\label{lesson:MR-C83-uchalne}



\begin{quote}
Before the new one: say the Marathi for to teach, then the Marathi for to make.
\end{quote}

\subsection*{You'll want to know: \mr{उचलणे}}
\textbf{\mr{उचलणे}} — \emph{uchalṇe} — \textquotedblleft{}to lift\textquotedblright{}.

\begin{cousinweb}
\textbf{\mr{उचलणे}} (\emph{uchalṇe}) is a tadbhava, a word inherited and reshaped by sound change. It comes from \emph{*uccalyati}, \textquotedblleft{}goes up\textquotedblright{}, a form CDIAL reconstructs from its descendants (CDIAL 1642). CDIAL builds it on Sanskrit \emph{uccalati}, \textquotedblleft{}springs up\textquotedblright{}.
\end{cousinweb}

\begin{grammarlens}[title={What you've built}]
One more word for this chapter.
\end{grammarlens}

\subsection*{Your turn}
\begin{itemize}
\raggedright
  \item \emph{Say it:} \emph{uchalṇe}
  \item \emph{Say it:} \emph{uchalṇe}, once more
  \item \emph{Say it:} say \emph{uchalṇe}
  \item \emph{Recall:} say the Marathi for to teach, then the Marathi for to make, then say \emph{uchalṇe} again
\end{itemize}

\begin{tcolorbox}[breakable,colback=teal!4,colframe=teal!35!black,title={Before you move on}]
What does \mr{उचलणे} mean? (\textquotedblleft{}To lift\textquotedblright{}.) Say it once more.
\end{tcolorbox}
\section[\texorpdfstring{\mr{ओढणे} (oḍhṇe)}{ओढणे (oḍhṇe)}]{\mr{ओढणे} (oḍhṇe) — to pull}

\label{lesson:MR-C83-odhne}



\begin{quote}
Before the new one: say the Marathi for to make, then the Marathi for to lift.
\end{quote}

\subsection*{You'll want to know: \mr{ओढणे}}
\textbf{\mr{ओढणे}} — \emph{oḍhṇe} — \textquotedblleft{}to pull\textquotedblright{}.

\begin{cousinweb}
\textbf{\mr{ओढणे}} (\emph{oḍhṇe}) is a tadbhava, a word inherited and reshaped by sound change. It comes from \emph{*avakaḍḍhati}, \textquotedblleft{}pulls away\textquotedblright{}, a form CDIAL reconstructs from its descendants (CDIAL 14245).
\end{cousinweb}

\begin{grammarlens}[title={What you've built}]
One more word for this chapter.
\end{grammarlens}

\subsection*{Your turn}
\begin{itemize}
\raggedright
  \item \emph{Say it:} \emph{oḍhṇe}
  \item \emph{Say it:} \emph{oḍhṇe}, once more
  \item \emph{Say it:} say \emph{oḍhṇe}
  \item \emph{Recall:} say the Marathi for to make, then the Marathi for to lift, then say \emph{oḍhṇe} again
\end{itemize}

\begin{tcolorbox}[breakable,colback=teal!4,colframe=teal!35!black,title={Before you move on}]
What does \mr{ओढणे} mean? (\textquotedblleft{}To pull\textquotedblright{}.) Say it once more.
\end{tcolorbox}
\section[\texorpdfstring{\mr{ढकलणे} (ḍhakalṇe)}{ढकलणे (ḍhakalṇe)}]{\mr{ढकलणे} (ḍhakalṇe) — to push}

\label{lesson:MR-C83-dhakalne}



\begin{quote}
Before the new one: say the Marathi for to lift, then the Marathi for to pull.
\end{quote}

\subsection*{You'll want to know: \mr{ढकलणे}}
\textbf{\mr{ढकलणे}} — \emph{ḍhakalṇe} — \textquotedblleft{}to push\textquotedblright{}.

\begin{cousinweb}
\textbf{\mr{ढकलणे}} (\emph{ḍhakalṇe}) belongs with Marathi \emph{ḍhakā}, \textquotedblleft{}push\textquotedblright{}, which CDIAL lists under a reconstructed \emph{*dhakk-}, \textquotedblleft{}push, strike\textquotedblright{} (CDIAL 6701).
\end{cousinweb}

\begin{grammarlens}[title={What you've built}]
One more word for this chapter.
\end{grammarlens}

\subsection*{Your turn}
\begin{itemize}
\raggedright
  \item \emph{Say it:} \emph{ḍhakalṇe}
  \item \emph{Say it:} \emph{ḍhakalṇe}, once more
  \item \emph{Say it:} say \emph{ḍhakalṇe}
  \item \emph{Recall:} say the Marathi for to lift, then the Marathi for to pull, then say \emph{ḍhakalṇe} again
\end{itemize}

\begin{tcolorbox}[breakable,colback=teal!4,colframe=teal!35!black,title={Before you move on}]
What does \mr{ढकलणे} mean? (\textquotedblleft{}To push\textquotedblright{}.) Say it once more.
\end{tcolorbox}
\section[\texorpdfstring{\mr{लपवणे} (lapavṇe)}{लपवणे (lapavṇe)}]{\mr{लपवणे} (lapavṇe) — to hide}

\label{lesson:MR-C83-lapavne}



\begin{quote}
Before the new one: say the Marathi for to pull, then the Marathi for to push.
\end{quote}

\subsection*{You'll want to know: \mr{लपवणे}}
\textbf{\mr{लपवणे}} — \emph{lapavṇe} — \textquotedblleft{}to hide\textquotedblright{}.

\begin{cousinweb}
The origin of \textbf{\mr{लपवणे}} (\emph{lapavṇe}) is not settled. Neither CDIAL nor DEDR has an entry for it or for \emph{lapṇẽ}, \textquotedblleft{}to hide\textquotedblright{}.
\end{cousinweb}

\begin{grammarlens}[title={What you've built}]
One more word for this chapter.
\end{grammarlens}

\subsection*{Your turn}
\begin{itemize}
\raggedright
  \item \emph{Say it:} \emph{lapavṇe}
  \item \emph{Say it:} \emph{lapavṇe}, once more
  \item \emph{Say it:} say \emph{lapavṇe}
  \item \emph{Recall:} say the Marathi for to pull, then the Marathi for to push, then say \emph{lapavṇe} again
\end{itemize}

\begin{tcolorbox}[breakable,colback=teal!4,colframe=teal!35!black,title={Before you move on}]
What does \mr{लपवणे} mean? (\textquotedblleft{}To hide\textquotedblright{}.) Say it once more.
\end{tcolorbox}
\section[\texorpdfstring{\mr{जिंकणे} (jiṅkṇe)}{जिंकणे (jiṅkṇe)}]{\mr{जिंकणे} (jiṅkṇe) — to win}

\label{lesson:MR-C83-jinkne}



\begin{quote}
Before the new one: say the Marathi for to push, then the Marathi for to hide.
\end{quote}

\subsection*{You'll want to know: \mr{जिंकणे}}
\textbf{\mr{जिंकणे}} — \emph{jiṅkṇe} — \textquotedblleft{}to win\textquotedblright{}.

\begin{cousinweb}
\textbf{\mr{जिंकणे}} (\emph{jiṅkṇe}) is a tadbhava, a word inherited from Sanskrit and reshaped by sound change. It comes from \emph{jita}, \textquotedblleft{}won, conquered\textquotedblright{} (CDIAL 5224). CDIAL lists Marathi \emph{jĩkṇẽ}, \textquotedblleft{}to conquer\textquotedblright{}, from the Middle Indo-Aryan form \emph{*jikka-}.
\end{cousinweb}

\begin{grammarlens}[title={What you've built}]
One more word for this chapter.
\end{grammarlens}

\subsection*{Your turn}
\begin{itemize}
\raggedright
  \item \emph{Say it:} \emph{jiṅkṇe}
  \item \emph{Say it:} \emph{jiṅkṇe}, once more
  \item \emph{Say it:} say \emph{jiṅkṇe}
  \item \emph{Recall:} say the Marathi for to push, then the Marathi for to hide, then say \emph{jiṅkṇe} again
\end{itemize}

\begin{tcolorbox}[breakable,colback=teal!4,colframe=teal!35!black,title={Before you move on}]
What does \mr{जिंकणे} mean? (\textquotedblleft{}To win\textquotedblright{}.) Say it once more.
\end{tcolorbox}
